\documentclass{article}

% Required packages
\usepackage{amssymb}
\usepackage{amsmath}
\usepackage{graphicx}
\usepackage{geometry}
\usepackage{tikz}
\usepackage{array}
\usepackage{booktabs}
\usepackage{enumitem}
\usepackage{listings}
\usepackage{xcolor}
\usepackage{fancyhdr}

% Set page geometry
\geometry{a4paper, margin=1in}

% Configure listings for Python
\lstset{
  language=Python,
  basicstyle=\ttfamily\footnotesize,
  numbers=left,
  numberstyle=\tiny\color{gray},
  frame=single,
  breaklines=true,
  breakatwhitespace=true,
  captionpos=b,
  tabsize=4,
  showspaces=false,
  showstringspaces=false,
  showtabs=false,
  commentstyle=\color{gray}\textit,
  keywordstyle=\color{blue}\bfseries,
  stringstyle=\color{red}
}

% For centering images and tables
\usepackage{float}

\begin{document}

\pagestyle{fancy}
\chead{DSC 288: Capstone (Spring 2026)}
\lhead{Week 3: Team Member Grades}
\rhead{R. Scotty Rogers}

\section*{Week 3 Tasks}

\begin{enumerate}

  \item \textbf{Solar Feature Investigation (Solar MHD Turbulence)}
  \begin{itemize}
    \item Evaluated the potential of adding a solar-level feature to represent different categories of MHD turbulence, based on insights from Billcliff (2026).
    \item Decided against adding this feature at this stage due to the additional complexity it would introduce without a clear payoff.
    \item Choose GOES X-ray flux as the sole solar-level proxy due to time constraints of course.
  \end{itemize}

  \item \textbf{Project Restructuring}
  \begin{itemize}
    \item Restructured the project repository to follow a clean, standardized layout inspired by other successful projects.
    \item Created top-level directories: \texttt{configs/}, \texttt{data/}, \texttt{docs/}, \texttt{models/}, \texttt{notebooks/}, \texttt{results/}, \texttt{scripts/}, \texttt{src/}, and \texttt{tests/}.
    \item Consolidated documentation, research notes, and milestones under \texttt{docs/}.
    \item Renamed scripts with numbered prefixes to indicate pipeline order.
    \item Moved source code into a proper \texttt{src/swmi/} package structure for better modularity and maintainability.
  \end{itemize}

  \item \textbf{Literature Review: Lag Features and Multi-Altitude Propagation}
  \begin{itemize}
    \item Reviewed Wing (2016) for its treatment of Mutual Information (MI), Conditional Mutual Information (CMI), and Transfer Entropy as nonlinear feature importance measures, which are relevant for understanding solar wind--magnetosphere relationships.
    \item Consulted Hyndman and Athanasopoulos (\textit{Forecasting: Principles and Practice}) and Brownlee (\textit{Deep Learning for Time Series Forecasting}) for general guidance on lag construction in multi-step causal chains.
    \item Reviewed lag construction techniques across the Sun $\rightarrow$ L1 $\rightarrow$ GEO $\rightarrow$ LEO $\rightarrow$ ground propagation chain with a focus on avoiding data leakage. The improper lag alignment in a multi-altitude causal chain can easily introduce future information into training features.
  \end{itemize}

  \item \textbf{Bug Fix: Newell Coupling NaN Inflation}
  \begin{itemize}
    \item Identified an issue with the Newell coupling lag features, which had an unexpectedly high null value rate of over 30\%.
    \item Discovered that a single NaN anywhere in a 30-minute rolling window was causing the entire window's lag value to be null.
    \item Implemented a fix to compute the lag using only valid values within the window, and added a companion \texttt{n\_points} quality column to track data density per window.
    \item This fix reduced the null rate to roughly 8 to 10\%, significantly improving the quality of the lag features.
  \end{itemize}

\end{enumerate}

\end{document}